\documentclass[10pt]{article}
\usepackage[T1]{fontenc}
\usepackage[utf8]{inputenc}
\usepackage[french]{babel}
\usepackage{amsmath,amssymb,mathrsfs,mathpazo}
\usepackage{geometry}
\geometry{a4paper,margin=2cm}
\usepackage{graphicx,array,tabularx,tabularray,colortbl,multirow,multicol,booktabs,enumitem,pifont,ulem,fancybox,diagbox,cancel}
\usepackage{pstricks,pstricks-add,pst-plot,pst-tree,pst-node,pst-text,pst-func,pst-math,pst-bspline,pst-3dplot,pst-eucl,pst-solides3d}
\let\pscalculate\undefined
\usepackage{pst-bezier}
\usepackage{tikz,pgfplots,tkz-tab}
\usetikzlibrary{arrows,arrows.meta,patterns,calc,positioning,trees,decorations.markings,shapes.geometric,shapes.misc,shapes.arrows,intersections,decorations.pathmorphing,decorations.pathreplacing,backgrounds,shadows,babel}
\pgfplotsset{compat=1.18}
\usepgfplotslibrary{fillbetween}
\usepackage[np]{numprint}
\usepackage{listings,makecell,pas-tableur,twoopt}
\let\barre\undefined
\usepackage[pstricks]{tablvar}
\definecolor{darkgreen}{RGB}{0,100,0}
\providecommand{\red}{\color{red}}
\providecommand{\blue}{\color{blue}}
\providecommand{\green}{\color{green}}
\providecommand{\black}{\color{black}}
\providecommand{\white}{\color{white}}
\providecommand{\gray}{\color{gray}}
\providecommand{\cyan}{\color{cyan}}
\providecommand{\magenta}{\color{magenta}}
\providecommand{\yellow}{\color{yellow}}
\providecommand{\violet}{\color{violet}}
\providecommand{\orange}{\color{orange}}
\providecommand{\euro}{\texteuro}
\providecommand{\EUR}{\texteuro}
\providecommand{\ieme}{\textsuperscript{e}}
\providecommand{\degres}{\ensuremath{^\circ}}
\providecommand{\degree}{\ensuremath{^\circ}}
\providecommand{\vv}[1]{\overrightarrow{#1}}
\providecommand{\Coord}[1]{#1}
\providecommand{\decosix}{}
\providecommand{\starredbullet}{\ensuremath{\bullet}}
\pagestyle{empty}
\setlength{\parindent}{0pt}
\usepackage{adjustbox}\DeclareUnicodeCharacter{2236}{:}\DeclareUnicodeCharacter{2190}{\ensuremath{\leftarrow}}
\providecommand{\R}{\mathbb{R}}
\providecommand{\N}{\mathbb{N}}
\providecommand{\Z}{\mathbb{Z}}
\providecommand{\Q}{\mathbb{Q}}
\providecommand{\D}{\mathbb{D}}
\providecommand{\C}{\mathbb{C}}
\providecommand{\e}{{\mathrm{e}}}
\providecommand{\vect}[1]{\overrightarrow{#1}}
\providecommand{\vectt}[1]{\overrightarrow{\mathrm{#1}}}
\providecommand{\vvt}[1]{\overrightarrow{\mathrm{#1}}}
\providecommand{\np}[2][]{{#2~#1}}
\providecommand{\vv}[1]{{\overrightarrow{#1}}}
\renewcommand{\d}{{\mathrm{d}}}
\providecommand{\degres}{{^\circ}}
\providecommand{\degre}{{^\circ}}
\providecommand{\blue}{\color{blue}}
\providecommand{\red}{\color{red}}
\providecommand{\enskip}{\quad}
\providecommand{\llbracket}{\lbrack\!\lbrack}
\providecommand{\rrbracket}{\rbrack\!\rbrack}
\providecommand{\euro}{{\text{€}}}
\providecommand{\no}{{\text{n°}}}
\providecommand{\interval}[2]{{[#1;#2]}}
\providecommand{\intervFO}[2]{{[#1;#2[}}
\providecommand{\dotfill}{\ldots}

\providecommand{\ds}{\displaystyle}
\providecommand{\textding}[1]{#1}
\providecommand{\Oij}{\(\mathrm{O};\vec{\imath},\vec{\jmath}\)}
\providecommand{\Oijk}{\(\mathrm{O};\vec{\imath},\vec{\jmath},\vec{k}\)}
\providecommand{\Ouv}{\(\mathrm{O};\vec{u},\vec{v}\)}
\providecommand{\intervFF}[2]{\ensuremath{[#1;#2]}}
\renewcommand{\euro}{\ifmmode\text{\texteuro}\else\texteuro\fi}
\DeclareUnicodeCharacter{00B0}{\ifmmode^{\circ}\else\textdegree\fi}
\DeclareUnicodeCharacter{20AC}{\ifmmode\text{\texteuro}\else\texteuro\fi}
% Énoncés : Éducation nationale. Transcription et corrigés : APMEP, Denis Vergès et contributeurs. Découpage et classement : CoopMaths, Sébastien Lozano, Éric Elter et contributeurs. Adaptation de la présentation : Maths Explorer.
% CC BY-SA : https://coopmaths.fr/www/about/
% https://coopmaths.fr/alea/static/e3c/2020/tex/e3c_2020_00_sujet18_4.tex
\begin{document}
\noindent\textit{APMEP / CoopMaths. CC BY-SA. Présentation adaptée par Maths Explorer.}\par\medskip
\medskip

On donne ci-dessous les représentations graphiques respectives $\mathcal{C}_f$ et $\mathcal{C}_g$ de deux fonctions $f$ et $g$ définies sur $\R$ l’ensemble des nombres réels.

\begin{center}
\psset{xunit=1cm,yunit=0.25cm,labelFontSize=\scriptstyle,showorigin=false,labelsep=0.1pt}
\begin{pspicture}(-8,-13)(5,35)
\multido{\n=-7+1}{12}{\psline[linewidth=0.75pt,linecolor=lightgray](\n,-12.5)(\n,32.6)}
\multido{\n=-12+1}{45}{\psline[linewidth=0.75pt,linecolor=lightgray](-7,\n)(4,\n)}
\psaxes[linewidth=0.95pt,Dy=2]{->}(0,0)(-7,-12.5)(4,33)
\uput[ur](0.2,0.5){O}
\def\Func{x x x 3 add mul 9 sub mul 1 sub}
\psplot[plotpoints=2000,linewidth=1.25pt,linecolor=blue]{-5.2}{3.2}{\Func}

\psplot[plotpoints=2000,linewidth=1.25pt,linecolor=red]{-2}{1.21}{x x 10 mul 8 add mul 8 add}
\uput[ul](-4.2,16){\blue $\mathcal{C}_f$}
\uput[ur](1.2,30){\red $\mathcal{C}_g$}

\end{pspicture}
\end{center}

\smallskip

\begin{enumerate}
\item La fonction $f$ est définie sur $\R$ par 

\[f(x) = x^3 + 3x^2 - 9x - 1.\]
On admet qu’elle est dérivable sur $\R$ et on note $f'$ [\cancel{désigne}] sa fonction dérivée.
	\begin{enumerate}
		\item Calculer $f'(x)$.
		\item Déterminer le signe de $f'(x)$ en fonction du réel $x$.

En déduire le tableau de variation de la fonction $f$.
		\item Déterminer une équation de la droite $T$ tangente à $\mathcal{C}_f$ au point d’abscisse $-1$.
	\end{enumerate}
\item La fonction $g$ est une fonction polynôme du second degré, il existe donc trois réels $a$, $b$ et $c$
tels que : $g(x) = ax^2 + bx + c$ pour tout réel $x$ . On note $\Delta$ son discriminant.
	\begin{enumerate}
		\item Déterminer, à l’aide du graphique, le signe de $a$ et le signe de $\Delta$.
		\item La fonction $g$ est définie, pour tout réel $x$, par $g(x) = 10x^2 + 8x + 8$.

Démontrer que les courbes $\mathcal{C}_f$ et $\mathcal{C}_g$  ont un point commun d’abscisse $-1$ et qu’en ce
point elles ont la même tangente.
	\end{enumerate}
\end{enumerate}

\end{document}
